% =====================================================================================
% Donde entra cada dato en la reformulacion — figura propia en TikZ.
%
% ES LA CONTRAFIGURA DE LA 1.2 (cap01/mlqem_flujo). Alli el trazo rojo discontinuo es el
% valor ruidoso medido entrando COMO ENTRADA DEL MODELO, que es lo que obliga a ejecutar
% el circuito antes de poder corregirlo. Aqui el mismo trazo rojo discontinuo entra
% SOLO en la aritmetica final, a la derecha de la linea de puntos que separa el antes del
% despues. Las dos figuras se leen juntas: es el argumento de que la reformulacion no
% rompe el paradigma pre-ejecucion.
%
% CODIGO DE COLOR, el mismo que en el resto de la memoria:
%   azulunir     -> el dato que el modelo SI ve antes de ejecutar
%   acentofig    -> lo que se le HACE (el modelo, la aritmetica)
%   rojoclasico  -> el valor medido, que solo existe DESPUES de ejecutar
%                   (equivale al #b3261e del SVG de la figura 1.2)
%
% ⚠️ La flecha de f-sombrero CRUZA la linea de puntos a proposito: la prediccion nace
%    antes de ejecutar y se usa despues. Lo que NO cruza —y es el punto de la figura— es
%    el trazo rojo, que se queda entero a la derecha.
%
% ⚠️ \shorthandoff{<>} es obligatorio: `babel` en español hace activos < y >, y eso rompe
%    la sintaxis de puntas de flecha de TikZ. Va tambien en el preambulo, pero se repite
%    aqui para que el fichero sea autocontenido si se compila suelto.
% =====================================================================================
\begingroup
\shorthandoff{<>}%
\begin{tikzpicture}[
    >={Stealth[length=2mm]},
    line width=0.6pt,
    font=\footnotesize,
    caja/.style   = {draw=azulunir, line width=0.8pt, rounded corners=2pt,
                     align=center, inner sep=4pt, minimum height=8mm, minimum width=21mm},
    opera/.style  = {draw=acentofig, fill=acentofig!8, line width=0.9pt,
                     rounded corners=2pt, align=center, inner sep=5pt,
                     minimum height=9mm},
    medido/.style = {draw=rojoclasico, line width=0.8pt, rounded corners=2pt,
                     align=center, inner sep=4pt, minimum height=8mm},
  ]

  % --- Columna 1: lo que el modelo ve, y es TODO lo que ve --------------------------
  \node[caja] (circ) at (0,0.85)  {Circuito\\transpilado};
  \node[caja] (tele) at (0,-0.85) {Telemetría\\del chip};

  % --- Columna 2: el modelo ---------------------------------------------------------
  \node[opera, minimum width=15mm] (modelo) at (3.3,0) {Modelo};

  \draw[->, azulunir] (circ.east) -- (modelo.west);
  \draw[->, azulunir] (tele.east) -- (modelo.west);

  % --- Columna 3: la aritmetica de la correccion ------------------------------------
  \node[opera, minimum width=12mm] (arit) at (8.0,0) {$\displaystyle \frac{r}{\hat f}$};

  % La prediccion cruza la frontera: nace antes de ejecutar y se consume despues.
  \draw[->, acentofig] (modelo.east) -- node[above, midway, text=black] {$\hat f$}
        (arit.west);

  \draw[->] (arit.east) -- ++(0.8,0)
        node[right, align=left, inner sep=2pt] {$\langle O\rangle$\\mitigado};

  % --- El valor medido: entra SOLO aqui, y desde el lado de «despues» ----------------
  \node[medido] (med) at (8.0,-2.0) {$\langle O\rangle_{\text{ruidoso}} = r$\\medido al ejecutar};
  \draw[->, rojoclasico, dashed, line width=0.9pt] (med.north) -- (arit.south);

  % --- La frontera: antes / despues de ejecutar --------------------------------------
  \draw[dotted, gray, line width=0.9pt] (6.1,-2.85) -- (6.1,1.75);
  \node[gray, above] at (3.0,1.6)  {antes de ejecutar};
  \node[gray, above] at (8.4,1.6)  {después de ejecutar};

\end{tikzpicture}
\endgroup
